\documentclass[10pt,a4paper]{amsart}
\usepackage{amsmath,amssymb,amsthm,graphicx}
\usepackage{verbatim}
\newenvironment{reference}{\comment}{\endcomment}
%\newenvironment{reference}{}{}
\newenvironment{slogan}{\comment}{\endcomment}
\newenvironment{history}{\comment}{\endcomment}
% For commutative diagrams we use Xy-pic
\usepackage[all]{xy}

% We use 2cell for 2-commutative diagrams.
\xyoption{2cell}
% Theorem environments.
%
\theoremstyle{plain}
\newtheorem{theorem}[subsection]{Theorem}
\newtheorem{proposition}[subsection]{Proposition}
\newtheorem{lemma}[subsection]{Lemma}

\theoremstyle{definition}
\newtheorem{definition}[subsection]{Definition}
\newtheorem{example}[subsection]{Example}
\newtheorem{exercise}[subsection]{Exercise}
\newtheorem{situation}[subsection]{Situation}

\theoremstyle{remark}
\newtheorem{remark}[subsection]{Remark}
\newtheorem{remarks}[subsection]{Remarks}

\numberwithin{equation}{subsection}

% Macros
%
\def\lim{\mathop{\mathrm{lim}}\nolimits}
\def\colim{\mathop{\mathrm{colim}}\nolimits}
\def\Spec{\mathop{\mathrm{Spec}}}
\def\Hom{\mathop{\mathrm{Hom}}\nolimits}
\def\Ext{\mathop{\mathrm{Ext}}\nolimits}
\def\SheafHom{\mathop{\mathcal{H}\!\mathit{om}}\nolimits}
\def\SheafExt{\mathop{\mathcal{E}\!\mathit{xt}}\nolimits}
\def\Sch{\mathit{Sch}}
\def\Mor{\mathop{\mathrm{Mor}}\nolimits}
\def\Ob{\mathop{\mathrm{Ob}}\nolimits}
\def\Sh{\mathop{\mathit{Sh}}\nolimits}
\def\NL{\mathop{N\!L}\nolimits}
\def\CH{\mathop{\mathrm{CH}}\nolimits}
\def\proetale{{pro\text{-}\acute{e}tale}}
\def\etale{{\acute{e}tale}}
\def\QCoh{\mathit{QCoh}}
\def\Ker{\mathop{\mathrm{Ker}}}
\def\Im{\mathop{\mathrm{Im}}}
\def\Coker{\mathop{\mathrm{Coker}}}
\def\Coim{\mathop{\mathrm{Coim}}}

% Boxtimes
%
\DeclareMathSymbol{\boxtimes}{\mathbin}{AMSa}{"02}

%
% Macros for moduli stacks/spaces
%
\def\QCohstack{\mathcal{QC}\!\mathit{oh}}
\def\Cohstack{\mathcal{C}\!\mathit{oh}}
\def\Spacesstack{\mathcal{S}\!\mathit{paces}}
\def\Quotfunctor{\mathrm{Quot}}
\def\Hilbfunctor{\mathrm{Hilb}}
\def\Curvesstack{\mathcal{C}\!\mathit{urves}}
\def\Polarizedstack{\mathcal{P}\!\mathit{olarized}}
\def\Complexesstack{\mathcal{C}\!\mathit{omplexes}}
% \Pic is the operator that assigns to X its picard group, usage \Pic(X)
% \Picardstack_{X/B} denotes the Picard stack of X over B
% \Picardfunctor_{X/B} denotes the Picard functor of X over B
\def\Pic{\mathop{\mathrm{Pic}}\nolimits}
\def\Picardstack{\mathcal{P}\!\mathit{ic}}
\def\Picardfunctor{\mathrm{Pic}}
\def\Deformationcategory{\mathcal{D}\!\mathit{ef}}
\usepackage{cleveref}
\crefname{originallemma}{original lemma}{original lemmas}
\Crefname{originallemma}{Original lemma}{Original lemmas}
\makeatletter
\newlabel{algebra-lemma-colimit-category-fp-algebras}{{05N9}{source}{}{sourcepointer-1}{}}
\newlabel{algebra-lemma-colimit-category-fp-algebras@cref}{{[originallemma][0][]05N9}{[1][0][]source}{}{}{}}
\makeatother
\crefname{originallemma}{original lemma}{original lemmas}
\Crefname{originallemma}{Original lemma}{Original lemmas}
\makeatletter
\newlabel{algebra-lemma-colimit-finite}{{07RG}{source}{}{sourcepointer-2}{}}
\newlabel{algebra-lemma-colimit-finite@cref}{{[originallemma][0][]07RG}{[1][0][]source}{}{}{}}
\makeatother
\crefname{originallemma}{original lemma}{original lemmas}
\Crefname{originallemma}{Original lemma}{Original lemmas}
\makeatletter
\newlabel{algebra-lemma-colimit-lci}{{0C33}{source}{}{sourcepointer-3}{}}
\newlabel{algebra-lemma-colimit-lci@cref}{{[originallemma][0][]0C33}{[1][0][]source}{}{}{}}
\makeatother
\crefname{originallemma}{original lemma}{original lemmas}
\Crefname{originallemma}{Original lemma}{Original lemmas}
\makeatletter
\newlabel{discriminant-lemma-locally-comes-from-universal-finite}{{0FL4}{source}{}{sourcepointer-4}{}}
\newlabel{discriminant-lemma-locally-comes-from-universal-finite@cref}{{[originallemma][0][]0FL4}{[1][0][]source}{}{}{}}
\makeatother
\crefname{originallemma}{original lemma}{original lemmas}
\Crefname{originallemma}{Original lemma}{Original lemmas}
\makeatletter
\newlabel{discriminant-lemma-universal-finite-syntomic-etale}{{0FL2}{source}{}{sourcepointer-5}{}}
\newlabel{discriminant-lemma-universal-finite-syntomic-etale@cref}{{[originallemma][0][]0FL2}{[1][0][]source}{}{}{}}
\makeatother
\crefname{originallemma}{original lemma}{original lemmas}
\Crefname{originallemma}{Original lemma}{Original lemmas}
\makeatletter
\newlabel{discriminant-lemma-universal-finite-syntomic-smooth}{{0FL1}{source}{}{sourcepointer-6}{}}
\newlabel{discriminant-lemma-universal-finite-syntomic-smooth@cref}{{[originallemma][0][]0FL1}{[1][0][]source}{}{}{}}
\makeatother
\crefname{originalproposition}{original proposition}{original propositions}
\Crefname{originalproposition}{Original proposition}{Original propositions}
\makeatletter
\newlabel{discriminant-proposition-tate-map}{{0FKD}{source}{}{sourcepointer-7}{}}
\newlabel{discriminant-proposition-tate-map@cref}{{[originalproposition][0][]0FKD}{[1][0][]source}{}{}{}}
\makeatother
\crefname{originalsection}{original section}{original sections}
\Crefname{originalsection}{Original section}{Original sections}
\makeatletter
\newlabel{discriminant-section-discriminant}{{0BVH}{source}{}{sourcepointer-8}{}}
\newlabel{discriminant-section-discriminant@cref}{{[originalsection][0][]0BVH}{[1][0][]source}{}{}{}}
\makeatother
\crefname{originalsection}{original section}{original sections}
\Crefname{originalsection}{Original section}{Original sections}
\makeatletter
\newlabel{discriminant-section-finite-morphisms}{{0FKW}{source}{}{sourcepointer-9}{}}
\newlabel{discriminant-section-finite-morphisms@cref}{{[originalsection][0][]0FKW}{[1][0][]source}{}{}{}}
\makeatother
\crefname{originallemma}{original lemma}{original lemmas}
\Crefname{originallemma}{Original lemma}{Original lemmas}
\makeatletter
\newlabel{lemma-Garel-upstairs}{{0FLA}{source}{}{sourcepointer-10}{}}
\newlabel{lemma-Garel-upstairs@cref}{{[originallemma][0][]0FLA}{[1][0][]source}{}{}{}}
\makeatother
\title{Garel's trace map on the differential forms of a finite syntomic morphism, constructed and shown unique}
\author{The Stacks Project Authors}
\date{stacks-project:derham.tex:local-snapshot-20261012; GNU FDL 1.2 or later}
\begin{document}
\maketitle
\noindent\textit{Editorial note.} All delivered mathematical text is the author's own, copied from the pinned Stacks Project chapter file derham.tex; no statement or proof step was rewritten, reordered or removed.\par\medskip
\noindent\textit{Editorial note.} Preamble provenance: the pinned chapter file opens with \textbackslash{}input\{preamble\} and carries no preamble text of its own, so the project's single Stacks-wide preamble staged at sources/stacks-derham/preamble.tex is excerpted verbatim as three spans, lines 8-12 (verbatim package and the project's own reference/slogan/history display environments), lines 14-18 (Xy-pic setup) and lines 154-217 (theorem declarations and mathematical macros). The wrapper class and the amsmath/amssymb/amsthm/graphicx base packages are supplied by the assembler. No macro was invented or redefined.\par\medskip
\noindent\textit{Editorial note.} Same-chapter pointers are located at their \textbackslash{}label lines inside the pinned chapter file; cross-chapter pointers are located in the project's own tag list staged in the snapshot (sources/stacks-chow/tags.txt), at the verbatim lines carrying their canonical tags.\par\medskip
\noindent\textit{Allowed standard prerequisites (editorial).}\begin{itemize}
\item The upstairs Garel map from the differentials of the source to the pullback of the differentials of the target twisted by the determinant of the cotangent complex. Same chapter derham.tex, label lemma-Garel-upstairs (invoked at proof line 4062). It supplies the canonical map \$c\textasciicircum{}p\_\{Y/X\} : \textbackslash{}Omega\textasciicircum{}p\_\{Y/\textbackslash{}mathbf\{Z\}\} \textbackslash{}to f\textasciicircum{}*\textbackslash{}Omega\textasciicircum{}p\_\{X/\textbackslash{}mathbf\{Z\}\} \textbackslash{}otimes \textbackslash{}det(\textbackslash{}NL\_\{Y/X\})\$ used in the locally Noetherian case.
\item The Tate map identifying the determinant of the cotangent complex with the dualizing sheaf, and the trace of a finite morphism. Discriminants, label discriminant-proposition-tate-map (invoked at proof line 4068), Section \textbackslash{}ref\{discriminant-section-discriminant\} (invoked at proof line 4055) and Section \textbackslash{}ref\{discriminant-section-finite-morphisms\} (invoked at proof lines 4080 and 4092). They give the isomorphism \$c\_\{Y/X\} : \textbackslash{}det(\textbackslash{}NL\_\{Y/X\}) \textbackslash{}to \textbackslash{}omega\_\{Y/X\}\$ sending \$\textbackslash{}delta\$ to \$\textbackslash{}tau\$, and the evaluation \$f\_*\textbackslash{}omega\_\{Y/X\} \textbackslash{}to \textbackslash{}mathcal\{O\}\_X\$ sending \$\textbackslash{}tau\_\{Y/X\}\$ to \$\textbackslash{}operatorname\{Trace\}\_f(1)\$, which is what turns the twisted map into the map of the lemma and gives property (1).
\item The local universal finite syntomic model, and its smooth and etale members. Discriminants, labels discriminant-lemma-locally-comes-from-universal-finite (invoked at proof line 4122), discriminant-lemma-universal-finite-syntomic-smooth (invoked at proof line 4139) and discriminant-lemma-universal-finite-syntomic-etale (invoked at proof line 4140). They allow the uniqueness statement to be checked on the universal families \$V\_d \textbackslash{}to U\_d\$, whose total spaces are smooth over the base with the morphism etale over a dense open.
\item Absolute Noetherian approximation. Algebra, labels algebra-lemma-colimit-category-fp-algebras (invoked at proof line 4156), algebra-lemma-colimit-lci (invoked at proof line 4161) and algebra-lemma-colimit-finite (invoked at proof line 4161). They write an affine test base as a directed limit of Noetherian affine schemes and descend the finite syntomic morphism to a stage, which is how the construction and its base change compatibility are extended beyond the locally Noetherian case.
\end{itemize}\medskip
\ifdefined\hypertarget\hypertarget{sourcepointer-1}{}\fi
\noindent\textit{Source reference.} Original lemma 05N9: sources/stacks-chow/tags.txt, source lines 7298--7298. Navigation pointer for \textbackslash{}ref\{algebra-lemma-colimit-category-fp-algebras\} used inside the counted span (line 4156); the argument that consumes it is retained in the span.\par\medskip
\ifdefined\hypertarget\hypertarget{sourcepointer-2}{}\fi
\noindent\textit{Source reference.} Original lemma 07RG: sources/stacks-chow/tags.txt, source lines 9866--9866. Navigation pointer for \textbackslash{}ref\{algebra-lemma-colimit-finite\} used inside the counted span (line 4162); the argument that consumes it is retained in the span.\par\medskip
\ifdefined\hypertarget\hypertarget{sourcepointer-3}{}\fi
\noindent\textit{Source reference.} Original lemma 0C33: sources/stacks-chow/tags.txt, source lines 15178--15178. Navigation pointer for \textbackslash{}ref\{algebra-lemma-colimit-lci\} used inside the counted span (line 4161); the argument that consumes it is retained in the span.\par\medskip
\ifdefined\hypertarget\hypertarget{sourcepointer-4}{}\fi
\noindent\textit{Source reference.} Original lemma 0FL4: sources/stacks-chow/tags.txt, source lines 19484--19484. Navigation pointer for \textbackslash{}ref\{discriminant-lemma-locally-comes-from-universal-finite\} used inside the counted span (line 4122); the argument that consumes it is retained in the span.\par\medskip
\ifdefined\hypertarget\hypertarget{sourcepointer-5}{}\fi
\noindent\textit{Source reference.} Original lemma 0FL2: sources/stacks-chow/tags.txt, source lines 19482--19482. Navigation pointer for \textbackslash{}ref\{discriminant-lemma-universal-finite-syntomic-etale\} used inside the counted span (line 4140); the argument that consumes it is retained in the span.\par\medskip
\ifdefined\hypertarget\hypertarget{sourcepointer-6}{}\fi
\noindent\textit{Source reference.} Original lemma 0FL1: sources/stacks-chow/tags.txt, source lines 19481--19481. Navigation pointer for \textbackslash{}ref\{discriminant-lemma-universal-finite-syntomic-smooth\} used inside the counted span (line 4139); the argument that consumes it is retained in the span.\par\medskip
\ifdefined\hypertarget\hypertarget{sourcepointer-7}{}\fi
\noindent\textit{Source reference.} Original proposition 0FKD: sources/stacks-chow/tags.txt, source lines 19458--19458. Navigation pointer for \textbackslash{}ref\{discriminant-proposition-tate-map\} used inside the counted span (line 4068); the argument that consumes it is retained in the span.\par\medskip
\ifdefined\hypertarget\hypertarget{sourcepointer-8}{}\fi
\noindent\textit{Source reference.} Original section 0BVH: sources/stacks-chow/tags.txt, source lines 14910--14910. Navigation pointer for \textbackslash{}ref\{discriminant-section-discriminant\} used inside the counted span (line 4055); the argument that consumes it is retained in the span.\par\medskip
\ifdefined\hypertarget\hypertarget{sourcepointer-9}{}\fi
\noindent\textit{Source reference.} Original section 0FKW: sources/stacks-chow/tags.txt, source lines 19476--19476. Navigation pointer for \textbackslash{}ref\{discriminant-section-finite-morphisms\} used inside the counted span (line 4080); the argument that consumes it is retained in the span.\par\medskip
\ifdefined\hypertarget\hypertarget{sourcepointer-10}{}\fi
\noindent\textit{Source reference.} Original lemma 0FLA: sources/stacks-derham/pinned-original-source.tex, source lines 3738--3738. Navigation pointer for \textbackslash{}ref\{lemma-Garel-upstairs\} used inside the counted span (line 4062); the argument that consumes it is retained in the span.\par\medskip
\begin{lemma}
\label{lemma-Garel}
There exists a unique rule that to every finite syntomic
morphism of schemes $f : Y \to X$ assigns $\mathcal{O}_X$-module maps
$$
\Theta^p_{Y/X} :
f_*\Omega^p_{Y/\mathbf{Z}}
\longrightarrow
\Omega^p_{X/\mathbf{Z}}
$$
satisfying the following properties
\begin{enumerate}
\item the composition with
$\Omega^p_{X/\mathbf{Z}} \otimes_{\mathcal{O}_X} f_*\mathcal{O}_Y
\to f_*\Omega^p_{Y/\mathbf{Z}}$ is equal to
$\text{id} \otimes \text{Trace}_f$
where $\text{Trace}_f : f_*\mathcal{O}_Y \to \mathcal{O}_X$
is the map from
Discriminants, Section \ref{discriminant-section-discriminant},
\item the rule is compatible with base change.
\end{enumerate}
\end{lemma}

\begin{proof}
First, assume that $X$ is locally Noetherian. By
Lemma \ref{lemma-Garel-upstairs} we have a canonical map
$$
c^p_{Y/X} : \Omega_{Y/\mathbf{Z}}^p
\longrightarrow
f^*\Omega_{X/\mathbf{Z}}^p \otimes_{\mathcal{O}_Y} \det(\NL_{Y/X})
$$
By Discriminants, Proposition \ref{discriminant-proposition-tate-map}
we have a canonical isomorphism
$$
c_{Y/X} : \det(\NL_{Y/X}) \to \omega_{Y/X}
$$
mapping $\delta(\NL_{Y/X})$ to $\tau_{Y/X}$. Combined these maps give
$$
c^p_{Y/X} \otimes c_{Y/X} :
\Omega_{Y/\mathbf{Z}}^p
\longrightarrow
f^*\Omega_{X/\mathbf{Z}}^p \otimes_{\mathcal{O}_Y} \omega_{Y/X}
$$
By Discriminants, Section \ref{discriminant-section-finite-morphisms}
this is the same thing as a map
$$
\Theta_{Y/X}^p :
f_*\Omega_{Y/\mathbf{Z}}^p
\longrightarrow
\Omega_{X/\mathbf{Z}}^p
$$
Recall that the relationship between $c^p_{Y/X} \otimes c_{Y/X}$
and $\Theta_{Y/X}^p$ uses the evaluation map
$f_*\omega_{Y/X} \to \mathcal{O}_X$
which sends $\tau_{Y/X}$ to $\text{Trace}_f(1)$, see
Discriminants, Section \ref{discriminant-section-finite-morphisms}.
Hence property (1) holds. Property (2) holds for base changes by
$X' \to X$ with $X'$ locally Noetherian because both $c^p_{Y/X}$ and
$c_{Y/X}$ are compatible with such base changes. For $f : Y \to X$
finite syntomic and $X$ locally Noetherian,
we will continue to denote $\Theta^p_{Y/X}$ the solution we've just found.

\medskip\noindent
Uniqueness. Suppose that we have a finite syntomic morphism
$f: Y \to X$ such that $X$ is smooth over $\Spec(\mathbf{Z})$
and $f$ is \'etale over a dense open of $X$. We claim that
in this case $\Theta^p_{Y/X}$ is uniquely determined by property (1).
Namely, consider the maps
$$
\Omega^p_{X/\mathbf{Z}} \otimes_{\mathcal{O}_X} f_*\mathcal{O}_Y \to
f_*\Omega^p_{Y/\mathbf{Z}} \to
\Omega^p_{X/\mathbf{Z}}
$$
The sheaf $\Omega^p_{X/\mathbf{Z}}$ is torsion free (by the assumed
smoothness), hence it suffices to check that the restriction of
$\Theta^p_{Y/X}$ is uniquely determined over the dense open over
which $f$ is \'etale, i.e., we may assume $f$ is \'etale.
However, if $f$ is \'etale, then
$f^*\Omega_{X/\mathbf{Z}} = \Omega_{Y/\mathbf{Z}}$
hence the first arrow in the displayed equation is an isomorphism.
Since we've pinned down the composition, this guarantees uniqueness.

\medskip\noindent
Let $f : Y \to X$ be a finite syntomic morphism of locally Noetherian schemes.
Let $x \in X$. By Discriminants, Lemma
\ref{discriminant-lemma-locally-comes-from-universal-finite}
we can find $d \geq 1$ and a commutative diagram
$$
\xymatrix{
Y \ar[d] &
V \ar[d] \ar[l] \ar[r] &
V_d \ar[d] \\
X &
U \ar[l] \ar[r] &
U_d
}
$$
such that $x \in U \subset X$ is open, $V = f^{-1}(U)$
and $V = U \times_{U_d} V_d$. Thus $\Theta^p_{Y/X}|_V$
is the pullback of the map $\Theta^p_{V_d/U_d}$.
However, by the discussion on uniqueness above and
Discriminants, Lemmas
\ref{discriminant-lemma-universal-finite-syntomic-smooth} and
\ref{discriminant-lemma-universal-finite-syntomic-etale}
the map $\Theta^p_{V_d/U_d}$ is uniquely determined
by the requirement (1). Hence uniqueness holds.

\medskip\noindent
At this point we know that we have existence and uniqueness
for all finite syntomic morphisms $Y \to X$ with $X$ locally Noetherian.
We could now give an argument similar to the proof of
Lemma \ref{lemma-Garel-upstairs} to extend to general $X$.
However, instead it is possible to directly use
absolute Noetherian approximation to finish the proof.
Namely, to construct $\Theta^p_{Y/X}$
it suffices to do so Zariski locally on $X$ (provided we also
show the uniqueness). Hence we may assume $X$ is affine (small
detail omitted). Then we can write $X = \lim_{i \in I} X_i$
as the limit over a directed set $I$ of Noetherian affine schemes.
By Algebra, Lemma \ref{algebra-lemma-colimit-category-fp-algebras}
we can find $0 \in I$ and a finitely
presented morphism of affines $f_0 : Y_0 \to X_0$ whose base change to
$X$ is $Y \to X$. After increasing $0$ we may assume $Y_0 \to X_0$
is finite and syntomic, see
Algebra, Lemma \ref{algebra-lemma-colimit-lci} and
\ref{algebra-lemma-colimit-finite}. For $i \geq 0$ also the
base change $f_i : Y_i = Y_0 \times_{X_0} X_i \to X_i$ is finite syntomic.
Then
$$
\Gamma(X, f_*\Omega^p_{Y/\mathbf{Z}}) =
\Gamma(Y, \Omega^p_{Y/\mathbf{Z}}) =
\colim_{i \geq 0} \Gamma(Y_i, \Omega^p_{Y_i/\mathbf{Z}}) =
\colim_{i \geq 0} \Gamma(X_i, f_{i, *}\Omega^p_{Y_i/\mathbf{Z}})
$$
Hence we can (and are forced to) define $\Theta^p_{Y/X}$ as the colimit
of the maps $\Theta^p_{Y_i/X_i}$. This map is compatible with any
cartesian diagram
$$
\xymatrix{
Y' \ar[r] \ar[d] & Y \ar[d] \\
X' \ar[r] & X
}
$$
with $X'$ affine as we know this for the case of Noetherian affine schemes
by the arguments given above (small detail omitted; hint: if we also
write $X' = \lim_{j \in J} X'_j$ then for every $i \in I$ there is a $j \in J$
and a morphism $X'_j \to X_i$ compatible with the morphism $X' \to X$).
This finishes the proof.
\end{proof}

\end{document}
